\textbf{Время:} 02:00

\textbf{Тезис:} в плосковолновой обработке модельная корреляция получается как средний интерференционный вес тех же пар в выбранной ячейке импульсов.

\textbf{На экране:} три блока схемы слева направо; затем фаза слева в формуле, вес в правом блоке схемы и обведённое среднее справа; закончить нижней строкой о тех же n парах и принятом приближении.

\textbf{Текст:} На предыдущем экране у нас появились точки испускания и импульсы. Здесь показано, как из них построить корреляцию.

В левом блоке берём конечные пионы и составляем пары. В среднем блоке группируем пары по разности импульсов: выбираем небольшой диапазон каждой из трёх компонент q. Для каждой пары сохраняем также разность координат и времён испускания.

Посмотрим на формулу фазы внизу слева. Это уже знакомая обменная фаза: разность четырёхимпульсов умножаем на разность четырёхкоординат и делим на аш с чертой на c. В правом блоке по этой фазе вычисляем вес: единица плюс косинус.

Обведённая формула показывает следующий шаг. Веса всех пар выбранной ячейки складываем и делим на число этих же пар n. Так получаем средний вес в этой ячейке. Повторяем для остальных ячеек — и строим зависимость корреляции от q.

В эксперименте координат испускания нет, поэтому там мы использовали пары одного события и смешанный фон. В модели известные координаты позволяют вычислить обменный вес непосредственно.

Нижняя строка задаёт границу этого примера: взаимодействием пионов после испускания пренебрегаем. Поэтому здесь используем именно выведенный ранее косинус.

\textbf{Переход:} Теперь посмотрим, как такая корреляция выглядит в моём текущем расчёте.

\textbf{Справка:} Lisa et al. (2005), §2.8; Балинов (2025), исходный слайд 2. Используются свободные симметризованные плоские волны. Для кулоновского и сильного взаимодействий требуется соответствующая относительная волновая функция. Здесь q — полная разность четырёхимпульсов, x = (ct,r), а q·Δx — скалярное произведение с метрикой (+,−,−,−); при импульсах в ГэВ и координатах в фм фазу делят на \(\hc\). Группировка производится по пространственным компонентам q в выбранной системе отсчёта. n — число тех же модельных пар, а не число независимо смешанных пар. Для координат нужны последние взаимодействия или рождение конечного пиона в распаде; общая поздняя поверхность вывода может включать свободный пролёт. Схема объясняет метод и не является независимой проверкой старого кода.
